% vim: ai sw=2 spell spelllang=cs

\begin{frame}{Limity za běhu}

  \begin{enumerate}
    \item Nevázané k souborovému systému
    \begin{itemize}
      \item Obecné limity systému

      \pause

      \item \code{long sysconf(int name)} (\header{unistd.h})
      \begin{itemize}
	\item \code{name}: jméno limitu -- \code{_SC_ARG_MAX},
	  \code{_SC_OPEN_MAX}, \code{_SC_PAGESIZE}, \dots
	\item Návratová hodnota: limit
      \end{itemize}
    \end{itemize}

    \pause

    \item Vázané k souborovému systému
    \begin{itemize}
      \item Často limity daného souborového systému pro danou položku
      \begin{itemize}
	\item Pro adresář (d)
	\item Pro soubor (f)
      \end{itemize}

      \pause

      \item \code{long pathconf(const char *path, int name)}
	(\header{unistd.h})
      \begin{itemize}
	\item \code{path}: pro který soubor
	\item \code{name}: jméno limitu -- \code{_PC_LINK_MAX} (f),
	  \code{_PC_NAME_MAX} (d), \code{_PC_PATH_MAX} (d), \dots
	\item Návratová hodnota: limit
      \end{itemize}
    \end{itemize}
  \end{enumerate}

\end{frame}

\begin{frame}{Úkol}

  \heading{Limity za běhu}
  \begin{enumerate}
    \item Zjistěte velikost stránky v systému
    \item Vypište
    \item Zjistěte maximální délku jména souboru a cesty
    \item Vypište
    \item Vyzkoušejte
  \end{enumerate}

\end{frame}

\begin{frame}[fragile]{Limity procesu}

  \begin{itemize}
    \item Každý proces má limity, co smí
    \begin{itemize}
      \item Počet otevřených souborů
      \item Velikost paměti, \dots
    \end{itemize}

    \pause

    \item Dva druhy limitů
    \begin{itemize}
      \item Hard: jen root může zvyšovat
      \item Soft: uživatel může měnit až do výše hard
    \end{itemize}

    \pause

    \item V shellu: \cmd{ulimit} (Demo)

    \pause

    \item \code{int getrlimit(int resource, struct rlimit *rlim)} a
      \code{int setrlimit(int resource, const struct rlimit *rlim)}
    \begin{itemize}
      \item \code{resource}: který limit -- \code{RLIMIT_CPU},
	\code{RLIMIT_STACK}, \dots
      \item \code{rlim}: na co nastavit/návratová hodnota
      \item Návratová hodnota: \code{0} -- OK, jinak chyba
    \end{itemize}
  \end{itemize}

  \begin{block}{\code{struct rlimit}}
  \begin{lstlisting}
rlim_t rlim_cur;  /* Soft limit */
rlim_t rlim_max;  /* Hard limit (ceiling for rlim_cur) */
  \end{lstlisting}
  \end{block}

\end{frame}

\begin{frame}{Úkol}

  \heading{Limity za běhu}
  \begin{enumerate}
    \item Vypište si do dvou sloupců všechny limity
    \begin{itemize}
      \item \code{Soft0  Hard0}
      \item \code{Soft1  Hard1}
      \item \dots
    \end{itemize}

    \item Iterujte přes limity od limitu číslo \code{0}
    \item Využijte \code{RLIM_NLIMITS} jako strop
    \item Vyzkoušejte
    \item Porovnejte výstup s \cmd{ulimit -a} a \cmd{ulimit -aH}
  \end{enumerate}

\end{frame}

\begin{frame}{Priorita procesu}

  \begin{itemize}
    \item Některé procesy jsou výpočetně náročné

    \pause

    \item Ale nevadí jim, když je někdo upozadí

    \pause

    \item Každý proces má prioritu
    \begin{itemize}
      \item Priorita = 20 + nice (z pohledu jádra, demo)

      \pause

      \item Nice: iniciálně \code{0}, \code{< 0} zvyšuje prioritu, \code{> 0}
	snižuje
    \end{itemize}

    \pause

    \item Změna nice: \code{int nice(int inc)} (\header{unistd.h})
    \begin{itemize}
      \item \code{inc}: o kolik
    \end{itemize}

    \pause

    \item Zjištění/nastavení: (\header{unistd.h}) \\
      \code{int getpriority(int which, id_t who)} a
      \code{int setpriority(int which, id_t who, int nice)}
    \begin{itemize}
      \item \pozor: pracují také s nice!
      \item \code{which}: čeho -- \code{PRIO_PROCESS}, \code{PRIO_PGRP},
	\code{PRIO_USER}
      \item \code{who}: kterého (procesu, skupiny procesů, uživatele)
      \item \code{nice}: hodnota nice
    \end{itemize}

    \pause

    \item Další speciální nastavení plánovače: \doc{man 7 sched}
    \begin{itemize}
      \item \code{sched_setaffinity}, \code{sched_setscheduler},
	\code{sched_setparam}, \dots
    \end{itemize}

  \end{itemize}

\end{frame}

\begin{frame}{Úkol}

  \heading{Priorita procesu}
  \begin{enumerate}
    \item Vytvořte si program, který cyklí
    \item Každou miliontou iteraci něco vypište
    \begin{itemize}
      \item Popř.\ číslo upravte podle rychlosti stroje
    \end{itemize}

    \item Před cyklem programu zvyšte nice (\code{nice( > 0)})
    \begin{itemize}
      \item Poběží s nižší prioritou
    \end{itemize}

    \item Program přeložte a spusťte čtyřikrát

    \item Nyní spustťe 4 instance \cmd{cat /dev/zero >/dev/null}
    \item Změnilo se chování vašeho programu?
    \item Co se stane, když nice měnit nebudete?
  \end{enumerate}

\end{frame}

